% TOP_PROBLEM: {"id": 3100085, "title": "Let f(p;n,k) = C(n,k) p^(k) (1-p)^(n-k)", "category_id": 2, "category": {"id": 2, "name": "combinatorics", "display_name": "Combinatorics", "description": "Counting problems, graph theory, discrete structures.", "slug": "combinatorics", "order_index": 2, "created_at": "2026-07-31T15:26:25.671Z"}, "status": "open", "problem_number": "AMR-030-0085", "set_id": 13}
% FIRST_POSTED: 2026-10-02
% ENTRY_KIND: statement_only
% UnsolvedMath ID 3100085; source catalogue code AMR-030-0085
% !TeX program = lualatex
% Definition and sources only; this file is not an LLM solution attempt.
\documentclass[11pt,a4paper]{article}
\usepackage[margin=25mm]{geometry}
\usepackage{fontspec}
\setmainfont{lmroman10-regular.otf}[BoldFont=lmroman10-bold.otf,
 ItalicFont=lmroman10-italic.otf,BoldItalicFont=lmroman10-bolditalic.otf]
\usepackage{amsmath,amssymb,mathtools}
\usepackage{unicode-math}
\setmathfont{latinmodern-math.otf}
\AtBeginDocument{\let\setminus\smallsetminus}
\usepackage{xurl,hyperref}
\hypersetup{colorlinks=true,urlcolor=blue}
\setcounter{secnumdepth}{0}
\setlength{\parindent}{0pt}
\setlength{\parskip}{5pt}
\setlength{\emergencystretch}{3em}
\begin{document}
\section*{Let f(p;n,k) = C(n,k) p\textasciicircum{}(k) (1-p)\textasciicircum{}(n-k)}\label{3100085}
{\small UnsolvedMath ID 3100085 · AMR-030-0085 · Combinatorics · AMR Open Problem Lists}
\subsection{Definitions and mathematical statement}
For $n\ge1$, $0<p<1$, and $0\le k\le n$, let
\[
              b_{n,p}(k)=\binom nk p^k(1-p)^{n-k}
\]
be the probability mass function of a binomial random variable.
Exclude the symmetric parameter $p=1/2$. Can there exist a single
choice of $n,p$ and four pairwise distinct integers
$k,\ell,k',\ell'\in\{0,\ldots,n\}$ such that
\[
           b_{n,p}(k)=b_{n,p}(\ell),
                  \qquad
           b_{n,p}(k')=b_{n,p}(\ell')\,?
\]
The two common values need not be equal to one another. The
question asks whether two disjoint equality pairs can occur
within the same nonsymmetric binomial distribution.

Writing $r=p/(1-p)>0$, an equality for $k<\ell$ is equivalent to
\[
                r^{\ell-k}=\frac{\binom nk}{\binom n\ell}.
\]
Thus both pairs must be compatible with the same positive
$r\ne1$; satisfying separate equalities with different parameters
does not answer the question. The excluded endpoints $p=0,1$
would create many trivial zero masses, and $p=1/2$ produces
the usual reflection equalities. Indices at $0$ or $n$ themselves
are permitted when $0<p<1$. Resolving the problem means either
producing one exact example or proving that all such simultaneous
equalities are impossible.
\subsection{Short English statement}
Can one nonsymmetric binomial probability mass function have two equality pairs involving four distinct indices?
\subsection{Sources}
\begin{itemize}
\item[S1] ULAM.AI and the UnsolvedMath Contributors, supplied problems.json, record 3100085 (AMR-030-0085), AMR Open Problem Lists. Catalogue identity and source transcription; CC BY 4.0.
\url{https://huggingface.co/datasets/ulamai/UnsolvedMath}.
\item[S2] Cooper - Combinatorial Problems I Like (2020 snapshot), Probability, source-order bullet 85. Original source indexed by AMR-030-0085.
\url{https://people.math.sc.edu/cooper/combprob.html}.
\end{itemize}
\end{document}
